% Einheit 1 von 4 – Geradengleichung aufstellen und lesen (bank/geraden/e1.jsonl, zusammenbau v0.1)
%% TODO Zeitmarke Einheit 1: Marken-Zeile nicht lesbar
%% TODO Zweigzeile Teil 1: was hier gelernt wird (Satz in Schülersprache aus dem Einheitstitel) – Einheit 1
\einheitenkopf[e1]{Einheit 1 von 4 · Geradengleichung aufstellen und lesen}
\zweigzeile{Abitur GK}
\setcounter{aufgabe}{9}

%% TODO Titel: Ich-kann-Satz fehlt in der Bank (Platzhalter: Kettenname) – Nr. 10
%% TODO Anweisung: ein Satz über der ersten Teilaufgabe fehlt in der Bank – Nr. 10
\begin{aufgabe}{„Gerade oder Strecke?“}
\begin{teile}
\teil Ein Förderband wird im Modell durch $\vec x = (5 | 1 | 0) + r \cdot (-3 | 2 | 1)$ mit $0 \le r \le 1$ beschrieben. Beschreibt die Gleichung eine Gerade oder eine Strecke? Kreuze an und nenne Anfang und Ende. \kreuz{Gerade} \\ \kreuz{Strecke}
\end{teile}
\end{aufgabe}

%% TODO Titel: Ich-kann-Satz fehlt in der Bank (Platzhalter: Kettenname) – Nr. 11
%% TODO Anweisung: ein Satz über der ersten Teilaufgabe fehlt in der Bank – Nr. 11
\begin{aufgabe}{Geradengleichung aufstellen und lesen}
%% TODO \rechenplatz{n}: Schrittzahl des Grundfalls fehlt in der Bank (nur bei fester Schreibform, 2.3 b)
\begin{teile}
\teil Gegeben ist $g\colon \vec x = (4 | -1 | 2) + t \cdot (3 | 0 | -2)$. Welcher Vektor gibt die Richtung von $g$ an? \kreuz{$(4 | -1 | 2)$} \\ \kreuz{$(3 | 0 | -2)$} \\ \kreuz{$(7 | -1 | 0)$}
\teil Stelle eine Gleichung der Geraden $g$ durch $A(2 | 1 | 3)$ und $B(4 | 4 | 2)$ auf.
\teil Ein Skilift führt geradlinig von der Talstation $T(2 | -4 | 1)$ zur Bergstation $B(-6 | 20 | 13)$; 1 LE entspricht 10 m. Gib eine Gleichung der Liftstrecke an und beschreibe, was sie im Sachzusammenhang darstellt. \hfill (Abitur 2023 GK)
\teil Gegeben sind $A(5 | -2 | 3)$ und $B(5 | 4 | 3)$. Begründe, dass die Gerade durch $A$ und $B$ parallel zur $x_2$-Achse verläuft. \hfill (Abitur 2021 LK)
\teil Gegeben ist $g\colon \vec x = (3 | 1 | 2) + r \cdot (1 | 3 | -2)$, $r \in \mathbb{R}$. Gib eine Gleichung einer Geraden $p$ an, die echt parallel zu $g$ verläuft, und eine Gleichung einer Geraden $q$, die $g$ senkrecht schneidet. \hfill (Abitur 2019 GK)
\teil Die Geraden $g$ und $h$ verlaufen parallel mit dem Richtungsvektor $(1 | 4 | 2)$; $P(2 | -1 | 3)$ liegt auf $g$, $Q(10 | 3 | -1)$ auf $h$. Der Punkt $T$ liegt auf der Strecke $PQ$ mit $|TQ| = 3 \cdot |PT|$. Gib eine Gleichung der Geraden an, die durch $T$ parallel zu $g$ und $h$ verläuft. \hfill (Abitur 2022 GK)
\teil Der Punkt $P(3 | 2 | 2)$ liegt in der Ebene $E\colon \vec x = (2 | -1 | 1) + r \cdot (1 | 0 | 2) + t \cdot (0 | 3 | -1)$. Die Gerade $g$ liegt in $E$, geht durch $P$ und hat keinen Schnittpunkt mit der $x_1x_2$-Ebene. Gib eine Gleichung von $g$ an. \hfill (Abitur 2021 GK)
\end{teile}
\end{aufgabe}

%% TODO Titel: Ich-kann-Satz fehlt in der Bank (Platzhalter: Kettenname) – Nr. 12
%% TODO Anweisung: ein Satz über der ersten Teilaufgabe fehlt in der Bank – Nr. 12
\begin{aufgabe}{Geradengleichung aufstellen und lesen – Fehler finden, Begründen, Anwendung, Darstellungswechsel}
\begin{teile}
\teil Finde den Fehler. Gesucht ist eine Gerade $p$, die echt parallel zu $g\colon \vec x = (1 | 4 | -2) + r \cdot (2 | 1 | 3)$ verläuft. Toms Lösung: $p\colon \vec x = (3 | 5 | 1) + r \cdot (2 | 1 | 3)$.
\teil Der Richtungsvektor einer Geraden hat die dritte Komponente null. Begründe ohne Rechnung, dass die Gerade die $x_1x_2$-Ebene nicht schneidet oder ganz in ihr liegt.
\teil Eine Drohne fliegt geradlinig und gleichmäßig in 30 s von $A(0 | 0 | 20)$ nach $B(120 | 90 | 50)$ (Angaben in m). Gib eine Gleichung der Flugstrecke an, in der der Parameter die Zeit in Sekunden ist, und bestimme die Position nach 10 s.
\teil Zeichne die Gerade $g\colon \vec x = (2 | 0 | 3) + t \cdot (0 | 2 | -1)$ in das Koordinatensystem: markiere den Stützpunkt, den Pfeil des Richtungsvektors und die Punkte für $t = 1$ und $t = 2$.

\begin{ksys3}\end{ksys3}
\end{teile}
\end{aufgabe}
